\subsection{UNIVERSAL CHARACTERIZATION OF TENSOR PRODUCTS OF ALGEBRAS}

PROPOSITION 5. Let \((\mathbf{E}_i)_{i\in I}\) be a finite family of A-algebras and, for each \(i\in I\) , let \(e_i\) be the unit element of \(\mathbf{E}_i\) . For each \(i\in I\) , let \(u_i: \mathbf{E}_i \to \mathbf{E} = \bigotimes_{i\in I} \mathbf{E}_i\) be the A-linear mapping defined by

\[
u _ {i} (x _ {i}) = \bigotimes_ {j} x _ {j} ^ {\prime} \quad \text { with } x _ {i} ^ {\prime} = x _ {i} \text { and } x _ {j} ^ {\prime} = e _ {j} \text { for } j \neq i.
\]

\begin{enumerate}
\def\labelenumi{(\roman{enumi})}
\item
  The \(u_{i}\) are A-algebra isomorphisms; further, for \(i \neq j\) , the elements \(u_{i}(x_{i})\) and \(u_{j}(x_{j})\) are permutable in \(\mathbf{E}\) for all \(x_{i} \in \mathbf{E}_{i}\) and \(x_{j} \in \mathbf{E}_{j}\) and \(\mathbf{E}\) is generated by the union of the subalgebras \(u_{i}(\mathbf{E}_{i})\) .
\item
  Let \(\mathbf{F}\) be an A-algebra and, for all \(i \in \mathbf{I}\) , let \(v_i: \mathbf{E}_i \to \mathbf{F}\) be an A-algebra homomorphism, where the \(v_i\) are such that, for \(i \neq j\) , \(v_i(x_i)\) and \(v_j(x_j)\) are permutable in \(\mathbf{F}\) for all \(x_i \in \mathbf{E}_i\) and \(x_j \in \mathbf{E}_j\) . Then there exists one and only one A-algebra homomorphism \(w: \mathbf{E} \to \mathbf{F}\) such that
\end{enumerate}

\[
v _ {i} = w \circ u _ {i} \quad \text { for all } i \in I.\tag{4}
\]

\begin{enumerate}
\def\labelenumi{(\roman{enumi})}
\tightlist
\item
  The mapping \(u_{i}\) is an algebra homomorphism by definition of the multiplication on \(\mathbf{E}\) . If \(i \neq j, x_{i} \in \mathbf{E}_{i}, x_{j} \in \mathbf{E}_{j}\) , then
\end{enumerate}

\[
u _ {i} (x _ {i}) = \bigotimes_ {k} x _ {k} ^ {\prime} \quad \text { with } x _ {i} ^ {\prime} = x _ {i}, x _ {k} ^ {\prime} = e _ {k} \text { for } k \neq i,
\]

\[
u _ {j} (x _ {j}) = \bigotimes_ {k} x _ {k} ^ {\prime \prime} \quad \text { with } x _ {j} ^ {\prime \prime} = x _ {j}, x _ {k} ^ {\prime \prime} = e _ {k} \text { for } k \neq j.
\]

Clearly \(x_{k}^{\prime}x_{k}^{\prime \prime} = x_{k}^{\prime \prime}x_{k}^{\prime}\) for all \(k\in \mathbf{I}\) and hence \(u_{i}(x_{i})\) and \(u_{j}(x_{j})\) commute in E by formula (1) (no. 1) defining the multiplication in E. The last assertion follows from the relation \(\bigotimes_{i}x_{i} = \prod_{i\in \mathbf{I}}u_{i}(x_{i})\) .

\begin{enumerate}
\def\labelenumi{(\roman{enumi})}
\setcounter{enumi}{1}
\tightlist
\item
  For each \(i \in I\) , let \(x_{i}\) be an element of \(\mathbf{E}_{i}\) . The product \(\prod_{i \in I} v_{i}(x_{i})\) is then defined in \(\mathbf{F}\) independently of any ordering on \(\mathbf{I}\) since the algebra \(\mathbf{F}\) is associative and the elements \(v_{i}(x_{i})\) are pairwise permutable. The mapping \((x_{i})_{i \in I} \to \prod_{i \in I} v_{i}(x_{i})\) of \(\prod_{i \in I} \mathbf{E}_{i}\) into \(\mathbf{F}\) is obviously A-multilinear and there therefore exists one and only one A-linear mapping \(w: \mathbf{E} \to \mathbf{F}\) such that
\end{enumerate}

\[
w \left(\bigotimes_ {i} x _ {i}\right) = \prod_ {i} v _ {i} \left(x _ {i}\right).\tag{5}
\]

Now, the desired A-algebra homomorphism \(w: \mathbf{E} \to \mathbf{F}\) must satisfy (5), which follows from (4) and the fact that \(\bigotimes_{i} x_i = \prod_{i \in I} u_i(x_i)\) . This proves the uniqueness of \(w\) ; it remains to show that the A-linear mapping \(w\) defined by (5) is an A-algebra homomorphism and satisfies (4). The fact that \(w\) satisfies (4) is obvious: it suffices to apply (5) to the case where \(x_j = e_j\) for \(j \neq i\) and we obtain \(w(u_i(x_i)) = v_i(x_i)\) . Finally, \(w\) is an algebra homomorphism, for

\[
\begin{array}{r l} w \left(\left(\bigotimes_ {i} x _ {i}\right) \left(\bigotimes_ {i} y _ {i}\right)\right) & = w \left(\bigotimes_ {i} (x _ {i} y _ {i})\right) = \prod_ {i} v _ {i} (x _ {i} y _ {i}) \\ & = \prod_ {i} (v _ {i} (x _ {i}) v _ {i} (y _ {i})) = \left(\prod_ {i} v _ {i} (x _ {i})\right). \left(\prod_ {i} v _ {i} (y _ {i})\right) \end{array}
\]

since \(v_{i}(x_{i})\) commutes with \(v_{j}(y_{j})\) for \(j \neq i\) ; hence

\[
w \left(\left(\bigotimes_ {i} x _ {i}\right) \left(\bigotimes_ {i} y _ {i}\right)\right) = w \left(\bigotimes_ {i} x _ {i}\right). w \left(\bigotimes_ {i} y _ {i}\right)
\]

which, by linearity, completes the proof.

The ordered pair consisting of E and the canonical mapping \(\phi: (x_i) \mapsto \bigotimes_i x_i\) of \(\prod_i E_i\) into E is a solution of the universal mapping problem (Set Theory, IV, § 3, no. 1) where \(\Sigma\) is the species of A-algebra structure, the morphisms being the A-algebra homomorphisms and the \(\alpha\) -mappings the mappings \(\prod_i u_i\) of \(\prod_i E_i\) into an A-algebra such that the \(u_i\) are A-algebra homomorphisms and \(u_i(x_i)\) and \(u_j(x_j)\) commute for \(i \neq j\) , for all \(x_i \in E_i\) and \(x_j \in E_j\) .

COROLLARY. Let \((\mathbf{E}_i)_{i\in I}, (\mathbf{F}_i)_{i\in I}\) be two finite families of A-algebras and, for all \(i\in I\) , let \(f_i: \mathbf{E}_i \to \mathbf{F}_i\) be an algebra homomorphism. If \(u_i: \mathbf{E}_i \to \bigotimes_{j\in I} \mathbf{E}_{j}, v_i: \mathbf{F}_i \to \bigotimes_{j\in I} \mathbf{F}_j\) are the canonical homomorphisms, the mapping \(f = \bigotimes_{i} f_i\) (cf.~no. 1) is the unique A-algebra homomorphism such that \(f \circ u_i = v_i \circ f_i\) for all \(i\in I\) .

It suffices to note that the homomorphisms \(g_{i} = v_{i} \circ f_{i}\) are such that \(g_{i}(x_{i}) = v_{i}(f_{i}(x_{i}))\) and \(g_{j}(x_{j}) = v_{j}(f_{j}(x_{j}))\) commute for \(i \neq j, x_{i} \in \mathbf{E}_{i}\) and \(x_{j} \in \mathbf{E}_{j}\) ; then apply Proposition 5.

When, in Proposition 5, the algebra F is assumed to be commutative, the hypothesis that \(v_{i}(x_{i})\) and \(v_{j}(x_{j})\) are permutable for \(i \neq j\) is automatically satisfied. Hence, when F is commutative, there is a canonical bijection

\[
\operatorname{Hom} _ {\mathrm{A-alg.}} \left(\bigotimes_ {i} \mathrm{E} _ {i}, \mathrm{F}\right)\rightarrow \prod_ {i} \operatorname{Hom} _ {\mathrm{A-alg.}} (\mathrm{E} _ {i}, \mathrm{F}),\tag{6}
\]

namely the one which associates with every homomorphism \(w\) of \(\bigotimes_{i} E_{i}\) into \(F\) the family of \(w \circ u_{i}\) .

Note that if \(\mathbf{E}\) is a commutative A-algebra, the ring structure of \(\mathbf{E} \otimes_{\mathbf{A}} \mathbf{F}\) is the same as that of \(\mathbf{F}_{(\mathbf{E})}\) ( \(\S 1\) , no. 5).

